% -----------------------------------------------------------
\section{Intercession}
% -----------------------------------------------------------
	\label{INTERCESSION}
	
% ----------------------
\subsection{See also}
% ----------------------

	\hyperref[ATONEMENT]{Atonement}, \hyperref[RECONCILIATION]{Reconciliation},

% ----------------------
\subsection{1828 American Dictionary of the English Language} % *1828
% ----------------------
	\label{INTERCESSION-1828}

	\begin{compactenum}
		\item The act of interceding; mediation; interposition between parties at variance, with a view to reconciliation; prayer or solicitation to one party in favor of another, sometimes against another.
	\end{compactenum}

% ----------------------
\subsection{Oxford English Dictionary} % *OED
% ----------------------
	\label{INTERCESSION-OED}

	\begin{compactitem}
		\item[I.1.a] The action of interceding or pleading on behalf of (rarely against) another; entreaty, solicitation, or prayer for another; mediation.
		\item[I.1.b.] \textit{spec}. in religious use: Intercessory prayer.
		\item[I.1.b.] Loosely used for a petition or pleading on one's own behalf. \textit{Obsolete}.
	\end{compactitem}

% ----------------------
\subsection{Scriptural Usage} % *SCRIPTURES
% ----------------------
	\label{INTERCESSION-SCRIPTURES}

	% -----
	\subsubsection{Old Testament} % *OT
	% -----
		\label{INTERCESSION-OT}
	
		% --
		\paragraph{KJV}
		% --
			\label{INTERCESSION-OT-KJV}
	
			\begin{tabular}{ll}
				Total occurrences in the OT & 4
			\end{tabular}
			
			\begin{compactdesc}
				\item[Isaiah 53:12] Therefore will I divide him a portion with the great, and he shall divide the spoil with the strong; because he hath poured out his soul unto death: and he was numbered with the transgressors; and he bare the sin of many, and made intercession for the transgressors.
				\item[Jeremiah 7:16] Therefore pray not thou for this people, neither lift up cry nor prayer for them, neither make intercession to me: for I will not hear thee.
				\item[Jeremiah 27:18] But if they be prophets, and if the word of the LORD be with them, let them now make intercession to the LORD of hosts, that the vessels which are left in the house of the LORD, and in the house of the king of Judah, and at Jerusalem, go not to Babylon.
				\item[Jeremiah 36:25] Nevertheless Elnathan and Delaiah and Gemariah had made intercession to the king that he would not burn the roll: but he would not hear them.
			\end{compactdesc}
			
		% --
		\paragraph{Hebrew Bible}
		% --
			\label{INTERCESSION-HEBREW}
		
			%
			\subparagraph{\texorpdfstring{\he{p*Aga`}}{}}
			%
				\label{PAGA} % whatever is in step bible without periods
				
				In Isaiah 53:12 the verb is in hiphil form; the appropriate definition seems to be \#2 below from HALOT. Note the parallel between this verb in Isaiah 53:6 and then again in 53:12 (meanings \#1 and \#2 below), and note with that the connection to \hyperref[DC19-16]{DC 19:16ff} and \hyperref[2NEPHI-2-9]{2 Nephi 2:9}.
			
				\begin{compactdesc}
					\item[HALOT 7456, PDF 2051] \textbf{}
						\begin{compactitem}
							\item[1] to let something hurt someone (Isaiah 53:6)
							\item[2] to look after someone, intercede for, with \he{b*:} and \he{'El}/\he{`E:l} (Isaiah 53:12)
							\item[3] to urge someone
						\end{compactitem}
					\item[BDB 6293, PDF 1935] \textbf{}
						\begin{compactitem}
							\item[1] cause to light upon (Isaiah 53:6) 
							\item[2] cause one to entreat
							\item[3] make entreaty
							\item[4] interpose (Isaiah 53:12)
						\end{compactitem}
					\item[DCH PDF 3806] \textbf{}
						\begin{compactitem}
							\item[1] cause to come upon, visit (Isaiah 53:6)
							\item[2] make entreaty (Isaiah 53:12)
							\item[3] cause to entreat
						\end{compactitem}
				\end{compactdesc}
				
				\begin{compactdesc}
					\item[Some OT uses] \textbf{}
						\begin{compactdesc}
							\item[Jeremiah 7:16] make intercession
							\item[Jeremiah 27:18] make intercession
							\item[Jeremiah 36:25] make intercession
						\end{compactdesc}
				\end{compactdesc}
			
		% % --
		% \paragraph{Septuagint}
		% % --
			% \label{INTERCESSION-LXX}
		
			% %
			% \subparagraph{\texorpdfstring{\gr{sample word}}{}}
			% %
		
	% -----
	\subsubsection{New Testament} % *NT
	% -----
		\label{INTERCESSION-NT}
	
		% --
		\paragraph{KJV}
		% --
			\label{INTERCESSION-NT-KJV}
	
			\begin{tabular}{ll}
				Total occurrences in the NT & 6
			\end{tabular}
			
			\begin{compactdesc}
				\item[Romans 8:26] Likewise the Spirit also helpeth our infirmities: for we know not what we should pray for as we ought: but the Spirit itself maketh intercession for us with groanings which cannot be uttered.
				\item[Romans 8:27] And he that searcheth the hearts knoweth what is the mind of the Spirit, because he maketh intercession for the saints according to the will of God.
				\item[Romans 8:34]  Who is he that condemneth?  It is Christ that died, yea rather, that is risen again, who is even at the right hand of God, who also maketh intercession for us.
				\item[Romans 11:2] God hath not cast away his people which he foreknew.  Wot ye not what the scripture saith of Elias?  how he maketh intercession to God against Israel, saying,
				\item[1 Timothy 2:1] I EXHORT therefore, that, first of all, supplications, prayers, intercessions, and giving of thanks, be made for all men;
				\item[Hebrews 7:25] Wherefore he is able also to save them to the uttermost that come unto God by him, seeing he ever liveth to make intercession for them.
			\end{compactdesc}
			
		% --
		\paragraph{Greek New Testament}
		% --
			\label{INTERCESSION-GREEK}
				
			%
			\subparagraph{\texorpdfstring{\gr{ὑπερεντυγχάνω}}{}}
			%
				\label{HUPERENTUNCHANO}
			
				\begin{compactdesc}
					\item[BDAG 919a] \textbf{}
						\begin{compactitem}
							\item to intercede in behalf of another, \textit{plead, intercede}
						\end{compactitem}
				\end{compactdesc}
				
				\begin{compactdesc}
					\item[Some NT uses] \textbf{}
						\begin{compactdesc}
							\item[Romans 8:26] \gr{Ὡσαύτως δὲ καὶ τὸ πνεῦμα συναντιλαμβάνεται τῇ ἀσθενείᾳ ἡμῶν· τὸ γὰρ τί προσευξώμεθα καθὸ δεῖ οὐκ οἴδαμεν, ἀλλὰ αὐτὸ τὸ πνεῦμα ὑπερεντυγχάνει στεναγμοῖς ἀλαλήτοις,}
						\end{compactdesc}
				\end{compactdesc}
				
			%
			\subparagraph{\texorpdfstring{\gr{ἐντυγχάνω}}{}}
			%
				\label{ENTUNCHANO}
			
				\begin{compactdesc}
					\item[BDAG 301b] \textbf{}
						\begin{compactitem}
							\item to make an earnest request through contact with the person approached, \textit{approach} or \textit{appeal to someone}
						\end{compactitem}
				\end{compactdesc}
				
				\begin{compactdesc}
					\item[Some NT uses] \textbf{}
						\begin{compactdesc}
							\item[Romans 8:27] \gr{ὁ δὲ ἐραυνῶν τὰς καρδίας οἶδεν τί τὸ φρόνημα τοῦ πνεύματος, ὅτι κατὰ θεὸν ἐντυγχάνει ὑπὲρ ἁγίων.} \textit{(BDAG says that this verse is spekaing about intercession by the Holy Spirit, but the opening of the verse is a masculine participle, so I don't see why it couldn't be talking about Christ here, expanding on his role as going beyond the Spirit's role, which clearly is spoken of in Romans 8:26. In addition, the verses which clearly do speak of Christ's intercession---like Romans 8:34 and Hebrews 7:25---use \gr{ἐντυγχάνω}, not \gr{ὑπερεντυγχάνω}, though in fairness they re clearly related, and the former construction uses \gr{ὑπὲρ}.)}
							\item[Romans 8:34] \gr{τίς ὁ κατακρινῶν; Χριστὸς ὁ ἀποθανών, μᾶλλον δὲ ἐγερθείς, ὅς καί ἐστιν ἐν δεξιᾷ τοῦ θεοῦ, ὃς καὶ ἐντυγχάνει ὑπὲρ ἡμῶν·}
							\item[Romans 11:2] \gr{οὐκ ἀπώσατο ὁ θεὸς τὸν λαὸν αὐτοῦ ὃν προέγνω. ἢ οὐκ οἴδατε ἐν Ἠλίᾳ τί λέγει ἡ γραφή, ὡς ἐντυγχάνει τῷ θεῷ κατὰ τοῦ Ἰσραήλ;} \textit{(This verse goes with others, like Jeremiah 27:18, where a prophet is spoken of as making intercession for a people.)}
							\item[Hebrews 7:25] \gr{ὅθεν καὶ σῴζειν εἰς τὸ παντελὲς δύναται τοὺς προσερχομένους δι᾽ αὐτοῦ τῷ θεῷ, πάντοτε ζῶν εἰς τὸ ἐντυγχάνειν ὑπὲρ αὐτῶν.} \textit{(Again of Christ's intercession. A beautiful verse; reminiscent of \hyperref[DC45-3]{DC 45:3--5}.)}
						\end{compactdesc}
				\end{compactdesc}
				
			%
			\subparagraph{\texorpdfstring{\gr{ἔντευξις}}{}}
			%
				\label{ENTEUXIS}
			
				\begin{compactdesc}
					\item[BDAG 300a] \textbf{}
						\begin{compactitem}
							\item[1] a formal request put to a high official or official body, \textit{petition, request}
							\item[2] \textit{prayer}
							\item[2.a] \textit{intercessory prayer}
							\item[2.b] generally, \textit{prayer}
							\item[2.c] \textit{prayer of thanksgiving}
							\item[2.d] \textit{power of intercession} (only in certain contexts)
						\end{compactitem}
				\end{compactdesc}
				
				\begin{compactdesc}
					\item[Some NT uses] \textbf{}
						\begin{compactdesc}
							\item[1 Timothy 2:1] \gr{Παρακαλῶ οὖν πρῶτον πάντων ποιεῖσθαι δεήσεις, προσευχάς, ἐντεύξεις, εὐχαριστίας, ὑπὲρ πάντων ἀνθρώπων,}
							\item[1 Timothy 4:5] \gr{ἁγιάζεται γὰρ διὰ λόγου θεοῦ καὶ ἐντεύξεως.}
						\end{compactdesc}
				\end{compactdesc}
		
	% -----
	\subsubsection{Book of Mormon} % *BOM
	% -----
		\label{INTERCESSION-BOM}
	
		\begin{tabular}{ll}
			Total occurrences in the BOM & 4
		\end{tabular}
		
		\begin{compactdesc}
			\item[2 Nephi 2:9] Wherefore, he is the firstfruits unto God, inasmuch as he shall make intercession for all the children of men; and they that believe in him shall be saved.
			\item[2 Nephi 2:10] And because of the intercession for all, all men come unto God; wherefore, they stand in the presence of him, to be judged of him according to the truth and holiness which is in him.  Wherefore, the ends of the law which the Holy One hath given, unto the inflicting of the punishment which is affixed, which punishment that is affixed is in opposition to that of the happiness which is affixed, to answer the ends of the atonement---
			\item[Mosiah 14:12] Therefore will I divide him a portion with the great, and he shall divide the spoil with the strong; because he hath poured out his soul unto death; and he was numbered with the transgressors; and he bore the sins of many, and made intercession for the transgressors. \textit{(This verse is taken from Isaiah 53:12. There are only slight differences---here, a ``;'' after ``death'' instead of a ``:'' as in Isaiah, and a bit later, Isaiah has ``and he bare the sin of many'' instead of ``and he bore the sins of many.'' This is slightly changed in Yale; see the \hyperref[MOSIAH-14-12-NOTE-VERSIONS]{note at Mosiah 14:12}.)}
			\item[Mosiah 15:8] And thus God breaketh the bands of death, having gained the victory over death; giving the Son power to make intercession for the children of men---
		\end{compactdesc}
		
	% -----
	\subsubsection{Doctrine and Covenants} % *DC
	% -----
		\label{INTERCESSION-DC}
	
		\begin{tabular}{ll}
			Total occurrences in the DC & 0
		\end{tabular}
		
	% -----
	\subsubsection{Pearl of Great Price} % *PGP
	% -----
		\label{INTERCESSION-PGP}
	
		\begin{tabular}{ll}
			Total occurrences in the PGP & 0
		\end{tabular}
		
% % ----------------------
% \subsection{Key Phrases} % *KEYPHRASES *PHRASES
% % ----------------------
	% \label{INTERCESSION-PHRASES}

	% \begin{compactdesc}
		% \item[]
	% \end{compactdesc}
	
% % ----------------------
% \subsection{Bible Dictionaries} % *DICTIONARIES
% % ----------------------
	% \label{INTERCESSION-BD}

% % ----------------------
% \subsection{Restoration Usage} % *RESTORATION
% % ----------------------
	% \label{INTERCESSION-RESTORATION}

	% % -----
	% \subsubsection{Joseph Smith} % JS
	% % -----
		% \label{INTERCESSION-JS}
	
		% \begin{compactdesc}
			% \item[DATE/SOURCE]
		% \end{compactdesc}
	
	% % -----
	% \subsubsection{Brigham Young} % BY
	% % -----
		% \label{INTERCESSION-BY}
	
		% \begin{compactdesc}
			% \item[DATE/SOURCE]
		% \end{compactdesc}
	
% ----------------------
\subsection{Commentary and Studies} % *COMMENTARY *STUDIES
% ----------------------
	\label{INTERCESSION-COMMENTARY}

	% -----
	\subsubsection{Key Points}
	% -----
		\label{INTERCESSION-KEYS}
	
		\begin{compactitem}
			\item The Son receives ``power to make intercession for the children of men'' (\hyperref[MOSIAH-15-8]{Mosiah 15:8}), apparently when the will of the Son is ``swallowed up in the will of the Father'' (\hyperref[MOSIAH-15-7]{Mosiah 15:7}).
				\begin{compactitem}
					\item In the will being swallowed up, God also thus ``breaketh the bands of death, having gained the victory over death'' (\hyperref[MOSIAH-15-8]{Mosiah 15:8}).
				\end{compactitem}
		\end{compactitem}